%% on Clauder erstellt und von mir überarbeitet
% ===================================================================
% Abschnitt "Daten, Häufigkeitsverteilungen und ihre Darstellung"
% ===================================================================

\newglossaryentry{def_merkmal_stat}{
    name={Was ist ein Merkmal im Kontext der beschreibenden Statistik?},
    description={Eine Eigenschaft, die an einer Versuchseinheit erhoben bzw.\ gemessen wird (z.B.\ Lage, Farbe, Alter, Größe). Ein Merkmal tritt in verschiedenen Ausprägungen auf.},
    type=dinge
}

\newglossaryentry{def_grundgesamtheit_stat_1}{
    name={Was ist die Grundgesamtheit?},
    description={Die Grundgesamtheit ist die Gesamtheit aller Versuchseinheiten, über die im Rahmen einer Datenerhebung eine Aussage getroffen werden soll.},
    type=dinge
}

\newglossaryentry{def_grundgesamtheit_stat_2}{
    name={Warum muss die Grundgesamtheit sorgfältig definiert werden?},
    description={Da sie je nach Merkmal und Fragestellung unterschiedlich abgegrenzt ist.},
    type=dinge
}

\newglossaryentry{def_stichprobe_stat}{
    name={Was ist eine Stichprobe, und warum wird meist sie statt der gesamten Grundgesamtheit ausgewertet?},
    description={Eine Stichprobe ist eine (in der Regel kleinere) Teilmenge von $n$ Versuchseinheiten aus der Grundgesamtheit.
    Man wertet sie aus, weil die vollständige Grundgesamtheit oft zu groß ist.
    Der Begriff "Stichprobe" wird üblicherweise sowohl für die gewählten Versuchseinheiten als auch für die zugehörigen Merkmalsausprägungen verwendet.},
    type=dinge
}

\newglossaryentry{def_urliste_stat}{
    name={Was ist eine Urliste?},
    description={Die Liste der ursprünglich (in Erhebungsreihenfolge, unsortiert) erfassten Daten einer Stichprobe, typischerweise dargestellt als Paare (Index/Versuchseinheit, Merkmalsausprägung).},
    type=dinge
}

\newglossaryentry{def_arten_von_merkmalen}{
    name={Welche Arten von Merkmalen unterscheidet man?},
    description={Quantitative Merkmale (z.B.\ Alter in Jahren oder Größe) und qualitative Merkmale (z.B.\ untergewichtig/normal/übergewichtig oder Farbe).},
    type=dinge
}

\newglossaryentry{def_arten_von_merkmalen_1}{
    name={Welche Arten von quantitative Merkmalen unterscheidet man?},
    description={Diskret (z.B.\ Alter in Jahren) und stetig (z.B.\ Größe). Qualitative Merkmale unterteilt man in
    \emph{ordinale} Merkmale/Rangmerkmale mit natürlicher Ordnung (z.B.\ untergewichtig/normal/übergewichtig) und \emph{nominale} Merkmale ohne natürliche Ordnung (z.B.\ Farbe).},
    type=dinge
}

\newglossaryentry{def_arten_von_merkmalen_2}{
    name={Welche Arten von qualitativen Merkmalen unterscheidet man?},
    description={Quantitative Merkmale unterteilt man in \emph{diskret} (z.B.\ Alter in Jahren) und \emph{stetig} (z.B.\ Größe). Qualitative Merkmale unterteilt man in \emph{ordinale} Merkmale/Rangmerkmale mit natürlicher Ordnung (z.B.\ untergewichtig/normal/übergewichtig) und \emph{nominale} Merkmale ohne natürliche Ordnung (z.B.\ Farbe).},
    type=dinge
}

\newglossaryentry{def_absolute_haeufigkeit}{
    name={Wie ist die absolute Häufigkeit $H_n(a_i)$ einer Ausprägung $a_i$ in einer Stichprobe $x_1,\dots,x_n$ definiert?},
    description={$H_n(a_i) = \#\{j\in\{1,\dots,n\}: x_j=a_i\}$ -- die Anzahl der Stichprobenwerte, die gleich $a_i$ sind.},
    type=dinge
}

\newglossaryentry{def_relative_haeufigkeit}{
    name={Wie ist die relative Häufigkeit $h_n(a_i)$ einer Ausprägung $a_i$ definiert?},
    description={$h_n(a_i) = \dfrac1n H_n(a_i)$.},
    type=dinge
}

\newglossaryentry{bem_summe_haeufigkeiten}{
    name={Welche Summeneigenschaften gelten für die absoluten bzw.\ relativen Häufigkeiten aller Ausprägungen $a_1,\dots,a_s$?},
    description={$H_n(a_1)+\dots+H_n(a_s) = n$ und $h_n(a_1)+\dots+h_n(a_s) = 1$.},
    type=dinge
}

\newglossaryentry{bem_wahl_der_darstellung}{
    name={Wann sind Kreis- bzw.\ Balkendiagramm für die Darstellung einer Häufigkeitsverteilung ungeeignet, und was verwendet man stattdessen?},
    description={Bei sehr vielen verschiedenen Ausprägungen (z.B.\ bei einem quasi-stetigen Merkmal wie Alter in Jahren) werden Kreis- und Balkendiagramme unübersichtlich. Stattdessen bildet man Klassen (Wertebereiche) und stellt die Häufigkeiten pro Klasse in einem Histogramm dar.},
    type=dinge
}

\newglossaryentry{bem_restklasse_bei_ausreissern}{
    name={Wie geht man bei der Bildung von Histogramm-Klassen mit Ausreißern in den Daten um?},
    description={Man bildet für die Ausreißer eine Restklasse. Für diese wird kein eigener Balken gezeichnet, sondern nur vermerkt, wie viele Elemente sie enthält.},
    type=dinge
}

% ===================================================================
% Abschnitt "Empirische Häufigkeitsverteilungen und ihre Lagemaße"
% ===================================================================

\newglossaryentry{def_lagemass}{
    name={Wie ist ein Lagemaß formal definiert?},
    description={Eine Funktion $l:\mathbb{R}^n\to\mathbb{R}$ heißt Lagemaß, wenn $l(x_1+a,\dots,x_n+a) = l(x_1,\dots,x_n)+a$ für alle $(x_1,\dots,x_n)\in\mathbb{R}^n$ und alle $a\in\mathbb{R}$ gilt (Translationsäquivarianz): Verschiebt man alle Stichprobenwerte um $a$, verschiebt sich auch der Wert des Lagemaßes um $a$.},
    type=dinge
}

\newglossaryentry{bem_welche_mittel_sind_lagemasse}{
    name={Welche der drei klassischen Mittelwerte (arithmetisches, geometrisches, harmonisches Mittel) sind Lagemaße im formalen Sinn?},
    description={Nur das arithmetische Mittel ist ein Lagemaß. Geometrisches Mittel und harmonisches Mittel sind KEINE Lagemaße -- sie erfüllen die Translationsäquivarianz-Eigenschaft $l(x+a)=l(x)+a$ im Allgemeinen nicht.},
    type=dinge
}

\newglossaryentry{satz_schwerpunkteigenschaft_arith_mittel}{
    name={Wie lautet die Schwerpunkteigenschaft (Nulleigenschaft) des arithmetischen Mittels?},
    description={$\sum_{i=1}^n (x_i-\bar x) = 0$ -- die Summe der Abweichungen aller Werte vom arithmetischen Mittel ist stets Null.},
    type=dinge
}

\newglossaryentry{satz_minimierungseigenschaft_arith_mittel}{
    name={Wie lautet die Minimierungseigenschaft des arithmetischen Mittels?},
    description={Das arithmetische Mittel minimiert die quadratische Abweichung der Stichprobenwerte von einer Konstanten: $\sum_{i=1}^n (x_i-\bar x)^2 = \inf_{c\in\mathbb{R}} \sum_{i=1}^n (x_i-c)^2$, d.h.\ $\sum_{i=1}^n(x_i-c)^2$ wird minimal für $c=\bar x$.},
    type=dinge
}

\newglossaryentry{def_modalwert}{
    name={Wie ist der Modalwert (Modus) einer Stichprobe definiert?},
    description={Eine Ausprägung $a$ heißt Modalwert/Modus des Merkmals $X$, wenn $H_n(a) \ge H_n(x_i)$ für alle $i=1,\dots,n$ gilt -- also wenn $a$ die Ausprägung mit der größten Häufigkeit in der Stichprobe ist.},
    type=dinge
}

\newglossaryentry{bem_eigenschaften_modalwert}{
    name={Welche Eigenschaften hat der Modalwert (Definitionsbereich, Eindeutigkeit, Lagemaß-Eigenschaft)?},
    description={Der Modalwert ist auch für nominale Merkmale definiert (er benötigt keine Ordnung). Er ist in der Regel nicht eindeutig bestimmt. Für quantitative Merkmale ist der Modalwert ein Lagemaß.},
    type=dinge
}

\newglossaryentry{def_unimodal}{
    name={Wann heißt eine Häufigkeitsverteilung unimodal?},
    description={Wenn der Modalwert $x_{\mathrm{mod}}$ eindeutig bestimmt ist, d.h.\ es gibt nur eine einzige Ausprägung mit maximaler Häufigkeit.},
    type=dinge
}

\newglossaryentry{def_median}{
    name={Wie ist der Median (Zentralwert) $x_{\mathrm{med}}$ einer Stichprobe $x_1,\dots,x_n$ definiert?},
    description={Mit der geordneten Stichprobe $x_{(1)}\le\dots\le x_{(n)}$: $x_{\mathrm{med}} = x_{\left(\frac{n+1}{2}\right)}$ falls $n$ ungerade, und $x_{\mathrm{med}} = \frac12\left[x_{(n/2)}+x_{(n/2+1)}\right]$ falls $n$ gerade.},
    type=dinge
}

\newglossaryentry{satz_median_ist_lagemass}{
    name={Ist der Median ein Lagemaß?},
    description={Ja, der Median ist ein Lagemaß.},
    type=dinge
}

\newglossaryentry{satz_median_haelfte_eigenschaft}{
    name={Welche "Halbierungseigenschaft" erfüllt der Median bezüglich der relativen Häufigkeiten?},
    description={$\sum_{k: a_k\le x_{\mathrm{med}}} h_n(a_k) \ge \tfrac12$ und $\sum_{k: a_k\ge x_{\mathrm{med}}} h_n(a_k) \ge \tfrac12$ -- mindestens die Hälfte der Stichprobenwerte ist kleiner-gleich, und mindestens die Hälfte ist größer-gleich dem Median.},
    type=dinge
}

\newglossaryentry{satz_median_charakterisierung}{
    name={Wie lässt sich der Median durch die Halbierungseigenschaft eindeutig charakterisieren (auch wenn $a$ selbst kein beobachteter Wert ist)?},
    description={Erfüllt ein $a\in\mathbb{R}$ die beiden Ungleichungen $\sum_{k:a_k\le a} h_n(a_k)\ge\tfrac12$ und $\sum_{k:a_k\ge a} h_n(a_k)\ge\tfrac12$, so folgt $x_{(\lceil n/2\rceil)} \le a \le x_{(\lceil n/2\rceil+1)}$. $a$ muss dabei keine tatsächlich beobachtete Ausprägung sein.},
    type=dinge
}

\newglossaryentry{satz_median_minimierungseigenschaft}{
    name={Welche Minimierungseigenschaft besitzt der Median?},
    description={Der Median minimiert die Summe der absoluten Abweichungen: $\sum_{i=1}^n |x_i-c| \to \min$ für $c$ im Intervall $\left[x_{\left(\lfloor\frac{n+1}{2}\rfloor\right)}, x_{\left(\lceil\frac{n+1}{2}\rceil\right)}\right]$.},
    type=dinge
}

% ===================================================================
% Abschnitt "Empirische Häufigkeitsverteilungen und ihre Streumaße"
% ===================================================================

\newglossaryentry{bem_wozu_streumasse}{
    name={Wozu benötigt man Streumaße zusätzlich zu Lagemaßen?},
    description={Lagemaße beschreiben nur die zentrale Tendenz einer Verteilung, sagen aber nichts darüber aus, wie stark die Daten um diesen zentralen Wert streuen. Streumaße erfassen diese Variabilität/Streubreite der Beobachtungswerte.},
    type=dinge
}

\newglossaryentry{def_streumass}{
    name={Wie ist ein Streumaß formal definiert?},
    description={Eine Funktion $s:\mathbb{R}^n\to\mathbb{R}$ heißt Streumaß, wenn $s(x_1+a,\dots,x_n+a) = s(x_1,\dots,x_n)$ für alle $(x_1,\dots,x_n)\in\mathbb{R}^n$ und alle $a\in\mathbb{R}$ gilt (Translationsinvarianz): Verschiebt man alle Stichprobenwerte um $a$, ändert sich der Wert des Streumaßes nicht.},
    type=dinge
}

\newglossaryentry{def_spannweite}{
    name={Wie ist die Spannweite einer Stichprobe definiert, und ist sie robust gegenüber Ausreißern?},
    description={Spannweite $= x_{\max}-x_{\min} = x_{(n)}-x_{(1)}$. Sie ist ein Streumaß, aber NICHT robust gegenüber Ausreißern.},
    type=dinge
}

\newglossaryentry{def_medianabweichung}{
    name={Wie ist die Medianabweichung definiert, und ist sie robust gegenüber Ausreißern?},
    description={Man bildet die Hilfsstichprobe $y_i=|x_i-x_{\mathrm{med}}|$; die Medianabweichung ist der Median $y_{\mathrm{med}}$ dieser Hilfsstichprobe ("Median der absoluten Abweichungen vom Median"). Sie ist robust gegenüber Ausreißern.},
    type=dinge
}

\newglossaryentry{def_mittlere_absolute_abweichung_vom_median}{
    name={Wie ist die mittlere absolute Abweichung vom Median definiert, und ist sie robust?},
    description={$\bar s = \dfrac1n\sum_{i=1}^n |x_i-x_{\mathrm{med}}|$ (arithmetisches Mittel der Hilfsstichprobe $y_i=|x_i-x_{\mathrm{med}}|$). Sie ist NICHT robust gegenüber Ausreißern -- obwohl der Median selbst robust ist.},
    type=dinge
}

\newglossaryentry{def_mittlere_absolute_abweichung_vom_mittel}{
    name={Wie ist die mittlere absolute Abweichung vom arithmetischen Mittel definiert, und ist sie robust?},
    description={$\dfrac1n\sum_{i=1}^n |x_i-\bar x|$. Sie ist NICHT robust gegenüber Ausreißern.},
    type=dinge
}

\newglossaryentry{def_stichprobenvarianz}{
    name={Wie ist die (Stichproben-)Varianz $s^2$ definiert, und ist sie robust gegenüber Ausreißern?},
    description={$s^2 = \dfrac1n\sum_{i=1}^n (x_i-\bar x)^2$ -- die mittlere quadratische Abweichung vom arithmetischen Mittel. Sie ist NICHT robust gegenüber Ausreißern.},
    type=dinge
}

\newglossaryentry{def_standardabweichung}{
    name={Wie ist die (Stichproben-)Standardabweichung definiert?},
    description={$s = \sqrt{s^2}$, die Wurzel der Stichprobenvarianz. Sie ist ebenfalls nicht robust gegenüber Ausreißern (da $s^2$ es nicht ist).},
    type=dinge
}

\newglossaryentry{satz_steinersche_formel}{
    name={Wie lautet die Steinersche Formel zur (oft einfacheren) Berechnung der Stichprobenvarianz?},
    description={$s^2 = \dfrac1n\sum_{i=1}^n x_i^2 - \bar x^2 = \dfrac1n\sum_{i=1}^n x_i^2 - \left(\dfrac1n\sum_{i=1}^n x_i\right)^2$.},
    type=dinge
}

\newglossaryentry{satz_streumasse_via_haeufigkeitsverteilung}{
    name={Wie lassen sich $\bar s$ (mittlere absolute Abweichung vom Median) und $s^2$ (Varianz) mithilfe der Häufigkeitsverteilung statt der Rohdaten berechnen?},
    description={$\bar s = \dfrac1n\sum_{k=1}^s |a_k-x_{\mathrm{med}}|\,H_n(a_k) = \sum_{k=1}^s |a_k-x_{\mathrm{med}}|\,h_n(a_k)$ und $s^2 = \dfrac1n\sum_{k=1}^s (a_k-\bar x)^2\,H_n(a_k) = \sum_{k=1}^s (a_k-\bar x)^2\,h_n(a_k)$ -- man summiert über die $s$ verschiedenen Ausprägungen statt über alle $n$ Beobachtungen.},
    type=dinge
}

\newglossaryentry{def_empirisches_p_quantil}{
    name={Wie ist das empirische $p$-Quantil $x_p$ definiert?},
    description={Für $p\in(0,1)$: $x_p = x_{(\lceil np\rceil)}$ falls $np\notin\mathbb{N}$, und $x_p = \dfrac12\left[x_{(np)}+x_{(np+1)}\right]$ falls $np\in\mathbb{N}$.},
    type=dinge
}

\newglossaryentry{bem_quantil_eigenschaften}{
    name={Wie hängt das $p$-Quantil mit dem Median zusammen, und welche Lagemaß-Eigenschaft hat es?},
    description={$x_{1/2} = x_{\mathrm{med}}$. Wie der Median ist auch das empirische $p$-Quantil ein robustes Lagemaß.},
    type=dinge
}

\newglossaryentry{satz_charakterisierung_p_quantil}{
    name={Wie lässt sich das $p$-Quantil analog zum Median charakterisieren?},
    description={Das $p$-Quantil erfüllt $\sum_{k:a_k\le x_p} h_n(a_k)\ge p$ und $\sum_{k:a_k\ge x_p} h_n(a_k)\ge 1-p$ -- mindestens $100p\%$ der Werte sind kleiner-gleich, mindestens $100(1-p)\%$ größer-gleich $x_p$. Umgekehrt: Erfüllt ein $a$ diese beiden Ungleichungen (für gegebenes $p$), so liegt $a$ zwischen $x_{(\lceil np\rceil)}$ und $x_{(\lceil np\rceil+1)}$; $a$ muss dabei keine beobachtete Ausprägung sein.},
    type=dinge
}

\newglossaryentry{def_p_quantilsabstand}{
    name={Was ist der $p$-Quantilsabstand, und wie heißt der Spezialfall $p=1/4$?},
    description={Für $p\in(0,\tfrac12)$ ist der $p$-Quantilsabstand definiert als $x_{1-p}-x_p$. Der Fall $p=1/4$ heißt Quartilsabstand.},
    type=dinge
}

\newglossaryentry{lem_p_quantilsabstand_ist_streumass}{
    name={Ist der $p$-Quantilsabstand ein Streumaß?},
    description={Ja, der $p$-Quantilsabstand ist ein Streumaß (erfüllt die Translationsinvarianz-Eigenschaft $s(x+a)=s(x)$).},
    type=dinge
}
